Este documento apresenta o memorial de cálculo completo para o dimensionamento de um pórtico rolante, elaborado como entrega final da disciplina de Projeto de Máquinas.

\section{Organização do documento}

O desenvolvimento do memorial está organizado conforme a sequência de dimensionamento do equipamento. Inicialmente, a seção de Dados de Entrada consolida as condições de projeto, hipóteses adotadas e parâmetros necessários para os cálculos. Em seguida, o Sistema de Elevação apresenta o dimensionamento dos componentes responsáveis pelo içamento da carga, incluindo cabos, polias, tambor, motor, freio, redutor, acoplamentos, eixo e elementos associados. A seção de Estrutura reúne as verificações dos elementos estruturais do pórtico, enquanto o Sistema de Translação trata dos componentes relacionados ao deslocamento do equipamento.

O projeto se baseia nas normas ABNT aplicáveis, reunidas nas referências bibliográficas, bem como em catálogos oficiais de fabricantes utilizados para a seleção e verificação de componentes comerciais. Os apêndices complementam o desenvolvimento principal: o Projeto Preliminar apresenta parâmetros predeterminados que servem como material auxiliar para algumas etapas do dimensionamento, e os demais apêndices fornecem dados de apoio empregados em verificações específicas, como propriedades de materiais, trilhos e rodas.

\section{Objetivos}
\begin{itemize}[leftmargin=1.25cm]
  \item Apresentar os critérios técnicos adotados.
  \item Demonstrar os cálculos de dimensionamento dos principais componentes.
  \item Consolidar verificações de segurança e desempenho.
\end{itemize}
